

\section{Le seuil de 98% (Nation & Laufer) : Impossibilité de lire au dictionnaire}
space{0.3cm}
oindent

L'enseignement du grec ancien, bien que fondé sur une tradition séculaire de rigueur philologique,
se heurte aujourd'hui à une problématique centrale qui menace sa pérennité académique et culturelle
: l'incapacité de la majorité des apprenants, même après plusieurs années d'étude, à lire des textes
authentiques avec fluidité et compréhension profonde sans le recours constant à des outils de
médiation externe. Cette recherche approfondie s'inscrit dans le cadre de la thèse portant sur la
didactique du grec ancien, et vise spécifiquement à étayer le point II.3.2, qui postule l'existence
d'un seuil critique de couverture lexicale – le seuil de 98\% identifié par Hu et Nation, affinant
les travaux antérieurs de Laufer – en deçà duquel la lecture autonome est cognitivement impossible.

L'hypothèse centrale examinée ici est que le modèle pédagogique traditionnel, souvent qualifié de
méthode Grammaire-Traduction, repose sur un malentendu fondamental concernant la nature cognitive de
la lecture. Ce modèle suppose implicitement que le manque de vocabulaire peut être compensé par
l'utilisation d'un dictionnaire et l'application analytique de règles grammaticales. Or, les données
issues de la linguistique appliquée, de la psychologie cognitive et des neurosciences démontrent que
cette approche mène à une surcharge cognitive systémique. Le recours au dictionnaire n'est pas une
simple aide ; c'est une interruption du processus de traitement de l'information qui empêche la
constitution du sens.

Ce rapport se propose de disséquer les mécanismes statistiques et cognitifs qui sous-tendent cette
affirmation. Il s'agira de démontrer, preuves à l'appui, que la lecture fluide exige un bagage
lexical massif et, surtout, \textit{automatisé} – relevant de la mémoire procédurale plutôt que
déclarative. Nous analyserons d'abord les fondements statistiques du seuil de 98\%, avant d'explorer
les coûts cognitifs de la consultation lexicale, pour enfin aborder les implications
neurobiologiques de l'automatisation et proposer des pistes pédagogiques concrètes pour combler le
fossé lexical qui sépare les manuels de débutants des textes de Platon ou de Thucydide.



\subsection{2\. la réalité statistique : analyse du seuil de 98\% (nation, laufer, hu)}

La quantification du vocabulaire nécessaire à la compréhension écrite a fait l'objet de recherches
rigoureuses en acquisition des langues secondes (SLA \- \textit{Second Language Acquisition}). Ces
études ont permis de passer d'intuitions pédagogiques à des modèles statistiques précis, définissant
la relation entre la couverture lexicale (le pourcentage de mots connus dans un texte) et la
compréhension.



\subsection{De 95\% à 98\% : l'évolution du consensus scientifique}

Les travaux pionniers de Batia Laufer, dès la fin des années 1980, ont établi les premières balises
de cette recherche. Dans son étude fondamentale de 1989, Laufer suggérait qu'un seuil de couverture
lexicale de 95\% était nécessaire pour une compréhension \og raisonnable\fg{} d'un texte
académique.\cite{II-3-2-ref1} Concrètement, cela signifie que sur 100 mots rencontrés, le lecteur
doit en connaître 95 immédiatement pour espérer déduire le sens des 5 mots inconnus par inférence
contextuelle. À ce niveau, le lecteur rencontre un mot inconnu tous les 20 mots, soit environ un mot
par phrase ou toutes les deux lignes. Bien que ce seuil permette une certaine forme de compréhension
globale, il exige encore un effort cognitif soutenu pour combler les lacunes.\cite{II-3-2-ref2}

Cependant, des recherches ultérieures, notamment celles menées par Hu et Nation (2000), ont réévalué
ce chiffre à la hausse, en particulier pour la lecture \og non assistée\fg{} ou la lecture plaisir,
qui se rapproche de l'objectif de fluidité visé par les humanistes pour le grec ancien. Leur étude,
intitulée \textit{Unknown Vocabulary Density and Reading Comprehension}, a démontré que pour lire un
texte de fiction sans aide extérieure et avec un niveau de compréhension adéquat, une couverture de
98\% est indispensable.\cite{II-3-2-ref4}

Le passage de 95\% à 98\% ne représente pas une augmentation triviale. À 98\%, le lecteur ne
rencontre qu'un mot inconnu tous les 50 mots. Cette densité permet de maintenir le fil narratif et
la structure syntaxique en mémoire de travail sans rupture. En revanche, à 95\%, les interruptions
sont deux fois et demie plus fréquentes. Hu et Nation concluent que si 95\% peut suffire pour une
lecture intensive assistée par un enseignant, 98\% est le seuil de l'autonomie
réelle.\cite{II-3-2-ref6} Schmitt, Jiang et Grabe (2011) ont confirmé ces résultats, montrant une
relation linéaire entre la couverture lexicale et la compréhension : chaque point de pourcentage
gagné entre 90\% et 100\% améliore directement la capacité du lecteur à saisir le sens du
texte.\cite{II-3-2-ref7}



\subsection{Implications statistiques pour le corpus grec ancien}

L'application de ces seuils au corpus du grec ancien révèle l'ampleur du défi pédagogique.
Contrairement aux langues modernes où un vocabulaire de base de 2 000 à 3 000 mots permet souvent
une communication efficace, la richesse lexicale et morphologique du grec ancien, couplée à la
diversité des genres (poésie épique, tragédie, histoire, philosophie), exige un volume de
vocabulaire beaucoup plus important.

Les analyses basées sur les corpus numériques (comme le \textit{Thesaurus Linguae Graecae} ou la
\textit{Perseus Digital Library}) indiquent les volumes suivants pour atteindre les seuils critiques
:

\begin{itemize}\item \textbf{Pour atteindre 95\% de couverture :} Il est estimé qu'un apprenant doit maîtriser entre 3 000 et 5 000 familles de mots pour lire des textes authentiques avec ce niveau de couverture.\cite{II-3-2-ref9}\item \textbf{Pour atteindre 98\% de couverture :} Le chiffre grimpe à environ 8 000 à 9 000 familles de mots.\cite{II-3-2-ref9}\end{itemize}Or, la réalité des cursus universitaires et scolaires actuels est bien en deçà de ces exigences. Un
étudiant ayant suivi deux années de grec ancien maîtrise typiquement entre 1 000 et 1 500 mots. Les
statistiques de couverture lexicale pour des auteurs comme Xénophon ou Platon montrent qu'avec 1 500
mots, la couverture textuelle oscille péniblement entre 70\% et 80\%.\cite{II-3-2-ref12}



\subsection{L'impossibilité statistique de l'inférence contextuelle}

Un argument pédagogique fréquent consiste à dire que les élèves doivent apprendre à \og deviner par
le contexte\fg{}. Cependant, les travaux de Laufer (1997) sur la \og détresse lexicale\fg{}
(\textit{lexical plight}) démontrent que cette stratégie est inopérante en dessous du seuil
critique. Pour deviner le sens d'un mot inconnu, le contexte immédiat (les mots environnants) doit
être transparent. Si le contexte lui-même contient des inconnues, la triangulation du sens devient
impossible.

L'étude de Liu et Nation (1985) et les travaux ultérieurs de Laufer (2020) ont quantifié cette
capacité d'inférence :

\begin{itemize}\item À \textbf{98\% de couverture}, le taux de succès de la devinette contextuelle est élevé.\item À \textbf{95\%}, il reste significatif mais demande plus d'effort.\item En dessous de \textbf{95\%}, et particulièrement autour de 90\% ou moins, le taux de réussite chute drastiquement. L'élève ne \og devine\fg{} plus ; il spécule au hasard.\cite{II-3-2-ref13}\end{itemize}Ainsi, demander à un étudiant qui possède une couverture de 80\% de \og lire\fg{} un texte de Platon
sans dictionnaire revient à lui demander de résoudre une équation où les variables inconnues sont
plus nombreuses que les constantes. Le recours au dictionnaire n'est donc pas une solution de
facilité, mais une nécessité absolue imposée par le déficit statistique. Cependant, comme nous le
verrons, cette nécessité technique est une catastrophe cognitive.

Si les statistiques de Nation et Laufer définissent le \textit{quoi} (le nombre de mots
nécessaires), la psychologie cognitive explique le \textit{pourquoi} (pourquoi le manque de
vocabulaire empêche la lecture). La lecture n'est pas une simple juxtaposition de mots ; c'est un
processus complexe de construction de sens qui repose lourdement sur la \textbf{Mémoire de Travail
(MdT)}.



\subsection{Théorie de la charge cognitive et effet de l'attention divisée}

La Théorie de la Charge Cognitive (\textit{Cognitive Load Theory} \- CLT), développée par John
Sweller, postule que la mémoire de travail humaine possède une capacité de traitement extrêmement
limitée. Elle ne peut gérer simultanément qu'un nombre restreint d'éléments (généralement entre 4 et
7) avant d'être saturée.\cite{II-3-2-ref15}

Dans une lecture fluide, la reconnaissance des mots est automatisée. Le cerveau récupère le sens et
la fonction grammaticale des mots depuis la mémoire à long terme sans effort conscient, libérant
ainsi la totalité des ressources de la mémoire de travail pour les processus de haut niveau :
l'analyse syntaxique, l'intégration sémantique, et la construction du modèle mental du texte (la
macro-structure).\cite{II-3-2-ref17}

L'utilisation d'un dictionnaire pendant la lecture provoque deux phénomènes délétères majeurs :

1. \textbf{L'Effet de l'Attention Divisée (\textit{Split-Attention Effect}) :} L'apprenant doit
scinder son attention entre deux sources d'information physiquement séparées : le texte source et
l'outil de référence (dictionnaire papier ou numérique). Sweller a démontré que cette nécessité
d'intégrer mentalement des sources disparates impose une charge cognitive intrinsèque très lourde.
L'esprit doit \og suspendre\fg{} le traitement du texte pour engager une nouvelle tâche de
recherche.\cite{II-3-2-ref19}

2. \textbf{Surcharge et Effondrement de la Mémoire de Travail :} Le processus de recherche lexicale
prend du temps. Or, les informations stockées dans la boucle phonologique de la mémoire de travail
sont volatiles ; elles s'effacent en quelques secondes sans répétition. Lorsqu'un étudiant s'arrête
sur un mot inconnu (ex: \textit{ἀπαλλαγήν}), cherche sa forme canonique, feuillette le dictionnaire,
trouve le sens, puis revient au texte, le début de la phrase grecque a déjà disparu de sa mémoire
vive. La structure syntaxique qu'il était en train de construire s'effondre. Il est contraint de
relire la phrase depuis le début, doublant ainsi l'effort cognitif.\cite{II-3-2-ref21}

C'est pourquoi la lecture avec dictionnaire est souvent décrite comme une expérience de \og
décodage\fg{} plutôt que de lecture. L'étudiant résout une série de micro-problèmes lexicaux isolés
mais échoue à saisir le sens global, car ses ressources cognitives sont entièrement consommées par
l'accès lexical de bas niveau.\cite{II-3-2-ref23}



\subsection{Rupture de la construction de la macro-structure}

Selon le modèle de Kintsch et van Dijk (1978), la compréhension de texte implique la construction
d'une \textbf{macro-structure} (le sens global, le fil narratif) à partir des micro-propositions (le
sens des phrases individuelles). Ce processus exige une fluidité dans l'apport
d'informations.\cite{II-3-2-ref25}

Lorsque la couverture lexicale est insuffisante (inférieure à 95-98\%), le flux d'information est
trop fragmenté pour permettre la mise à jour constante de la macro-structure. L'étudiant, interrompu
à chaque ligne, perd le fil de l'argumentation ou de la narration. Il ne \og lit\fg{} pas une
histoire ; il traduit une liste de mots. Cette absence de macro-structure cohérente empêche à son
tour l'inférence et la prédiction, qui sont des mécanismes essentiels de la lecture fluide. Le
lecteur ne peut pas anticiper ce qui va suivre car il n'a pas compris ce qui précède, enfermé dans
une myopie lexicale.\cite{II-3-2-ref25}



\subsection{Le coût cognitif spécifique de la morphologie grecque}

Le grec ancien ajoute une couche de complexité supplémentaire par rapport à des langues comme
l'anglais ou l'espagnol : sa morphologie flexionnelle riche. Pour chercher un mot dans le
dictionnaire, l'étudiant doit d'abord effectuer une \textbf{lemmatisation} mentale. Face à une forme
comme \textit{ἠνέσχετο}, il doit identifier l'augment, le préfixe, la racine verbale, et la
désinence pour remonter au lemme \textit{ἀνέχω}.

Cette opération de \og rétro-ingénierie\fg{} morphologique est extrêmement coûteuse cognitivement.
Elle requiert une analyse explicite des règles grammaticales stockées en mémoire déclarative. Si
cette analyse n'est pas automatisée, elle sature immédiatement la mémoire de travail avant même que
l'étudiant n'ait ouvert son dictionnaire.\cite{II-3-2-ref28} Ainsi, le seuil de 98\% en grec
implique non seulement de connaître le sens des mots, mais de reconnaître instantanément leurs
formes fléchies sans analyse consciente.

Les preuves statistiques et cognitives sont corroborées par les découvertes en neurosciences de
l'éducation, notamment les travaux sur le \og cerveau lecteur\fg{} et les systèmes de mémoire.



\subsection{L'hypothèse du recyclage neuronal et la vwfa (dehaene)}

Stanislas Dehaene a mis en évidence le rôle central de la \textbf{Visual Word Form Area (VWFA)}, ou
\og boîte aux lettres du cerveau\fg{}, située dans le cortex occipito-temporal gauche. Cette zone,
issue du recyclage neuronal, se spécialise dans la reconnaissance visuelle des
mots.\cite{II-3-2-ref30}

Pour un lecteur expert, la reconnaissance d'un mot connu via la VWFA est ultra-rapide (moins de 200
millisecondes) et parallèle (toutes les lettres sont traitées simultanément). Cette voie de lecture,
dite \textbf{voie ventrale} ou directe, permet d'accéder au sens sans passer par une conversion
phonologique laborieuse. C'est le substrat neurologique de l'automaticité.\cite{II-3-2-ref32}

En revanche, face à un mot inconnu ou mal maîtrisé, le cerveau doit emprunter la \textbf{voie
dorsale} (pariéto-temporale), qui procède à un décodage séquentiel (lettre par lettre, syllabe par
syllabe) et à une conversion graphème-phonème. Cette voie est lente, attentionnelle et coûteuse en
énergie.\cite{II-3-2-ref34}

L'implication pour le grec ancien est cruciale : \og lire avec un dictionnaire\fg{} maintient le
cerveau de l'apprenant prisonnier de la voie dorsale. L'absence de reconnaissance instantanée
empêche l'activation efficace de la VWFA pour le script grec. Pour que la lecture devienne fluide,
il faut entraîner la VWFA par une exposition massive et répétée aux formes visuelles des mots, ce
que la méthode de traduction mot à mot ne permet pas.\cite{II-3-2-ref36}



\subsection{Le modèle déclaratif/procédural (ullman)}

Le modèle DP (\textit{Declarative/Procedural}) de Michael Ullman offre un cadre théorique pour
comprendre la nature de la compétence lexicale requise.

\begin{itemize}\item \textbf{Mémoire Déclarative :} Elle stocke les faits et les connaissances explicites (ex: \og Je sais que \textit{luō} signifie délier\fg{}). C'est une mémoire consciente, rapide à apprendre mais lente à récupérer en temps réel.\item \textbf{Mémoire Procédurale :} Elle gère les habiletés motrices et cognitives automatisées, y compris la grammaire et les séquences fréquentes (ex: comprendre \textit{elusamēn} sans analyser ses parties). C'est une mémoire implicite, lente à acquérir (nécessitant beaucoup de répétition) mais ultra-rapide à l'usage.\cite{II-3-2-ref38}\end{itemize}L'enseignement traditionnel du grec, basé sur l'apprentissage par cœur de tableaux de déclinaisons
et de listes de vocabulaire, sollicite presque exclusivement la mémoire déclarative. L'élève \og
sait\fg{} ses règles, mais il ne les a pas \og procéduralisées\fg{}. Lors de la lecture, il doit
récupérer consciemment ces règles, ce qui est trop lent pour la fluidité. Le seuil de 98\% de Nation
et Laufer correspond, en termes neurobiologiques, au moment où la majorité du traitement lexical et
morphologique a basculé vers le système procédural et la voie ventrale de
lecture.\cite{II-3-2-ref40}



\subsection{Automatisation de la morphologie flexionnelle}

Les études sur le traitement de la morphologie en L2 montrent que les apprenants tardifs ont
tendance à traiter les mots complexes (comme les verbes grecs) en les décomposant systématiquement,
là où les natifs les traitent souvent comme des formes holistiques.\cite{II-3-2-ref42} Pour
atteindre le seuil de fluidité, l'apprenant de grec doit automatiser non seulement le sens des
racines, mais aussi le \textit{parsing} morphologique. Il doit voir \textit{labōn} et penser \og
ayant pris\fg{} instantanément, sans passer par l'étape intermédiaire \og participe aoriste actif de
\textit{lambanō}\fg{}. Cette automatisation ne peut provenir que d'un input massif et
structuré.\cite{II-3-2-ref29}

L'analyse des statistiques lexicales révèle un fossé structurel – un véritable \og goulot
d'étranglement\fg{} – dans l'enseignement du grec ancien.



\subsection{Les chiffres du déficit}

Les manuels d'introduction standard (type \textit{Hermaion}, \textit{Reading Greek}, ou
\textit{Athenaze}) couvrent généralement les 1 000 à 1 500 mots les plus fréquents.

\begin{itemize}\item La liste \textbf{DCC Core Vocabulary} (Dickinson College) contient les 500 mots les plus fréquents, couvrant environ 65-70\% d'un texte littéraire standard.\cite{II-3-2-ref45}\item Les listes de fréquence établies par \textbf{Major (2013)} montrent que même avec 1 100 mots (visant 80\% de couverture), le lecteur reste confronté à un mot inconnu tous les cinq mots.\cite{II-3-2-ref47}\end{itemize}Pour passer de ce niveau (fin de licence ou de scolarité classique) au seuil de 95-98\% requis pour
lire Platon, Xénophon ou le Nouveau Testament sans dictionnaire, l'apprenant doit maîtriser entre
\textbf{3 000 et 9 000 familles de mots}.\cite{II-3-2-ref9} Il existe donc un \og no man's land\fg{}
pédagogique entre le niveau intermédiaire (1 500 mots) et le niveau de lecture autonome (8 000+
mots). C'est dans ce fossé que la majorité des étudiants abandonnent, découragés par la lourdeur de
la tâche de traduction qui ne s'allège jamais vraiment.\cite{II-3-2-ref48}



\subsection{Comparaison avec les langues vivantes}

Dans l'enseignement des langues vivantes (anglais, espagnol), ce fossé est comblé par une abondance
de \textit{Graded Readers} (lectures graduées) qui permettent à l'élève de passer progressivement de
500 à 1000, puis 2000, 3000 et 5000 mots. En grec ancien, ces ressources intermédiaires sont
historiquement rares, bien que la situation évolue. L'absence de ce \og pont\fg{} lexical oblige les
étudiants à sauter directement du manuel de base aux textes littéraires complexes, ce qui équivaut à
demander à un étudiant d'anglais langue étrangère de passer d'une méthode de niveau A2 directement à
Shakespeare.\cite{II-3-2-ref50}

Face à l'impossibilité statistique et cognitive de la méthode traditionnelle pour la lecture fluide,
la seule réponse viable est une restructuration pédagogique visant l'acquisition massive de
vocabulaire par le biais de l'\textbf{Input Compréhensible} (\textit{Comprehensible Input} \- CI).



\subsection{La lecture extensive et les lectures graduées}

La solution la plus directe pour atteindre le seuil de 98\% est d'augmenter le volume de lecture à
des niveaux de difficulté contrôlés.

\begin{itemize}\item \textbf{Lectures Graduées (\textit{Graded Readers}) :} Il est impératif de développer et d'utiliser des textes construits spécifiquement pour limiter la densité de vocabulaire inconnu. Des ouvrages récents tentent de combler ce vide (ex : les novellas de \textit{Comprehensible Classics}, les adaptations de \textit{Pugio Bruti} en grec, ou les lecteurs bibliques gradués).\cite{II-3-2-ref52} Ces textes permettent de consolider les mots fréquents par la répétition en contexte (incidental learning).\item \textbf{Lectures \og Tiered\fg{} ou Enchâssées (\textit{Embedded Readings}) :} Cette technique, popularisée par Laurie Clarcq et adaptée aux langues anciennes par des pédagogues comme Justin Slocum Bailey ou le projet \textit{Comprehensible Antiquity}, consiste à proposer plusieurs versions d'un même texte authentique.\item \textit{Niveau 1 :} Version très simplifiée, phrases courtes, vocabulaire restreint.\item \textit{Niveau 2 :} Introduction de subordination, vocabulaire plus riche.\item Niveau 3 : Texte original.\end{itemize}En lisant les versions préparatoires, l'élève construit le modèle mental et pré-active le
vocabulaire, ce qui fait monter artificiellement son taux de couverture lorsqu'il aborde le texte
original, rendant la lecture \og authentique\fg{} accessible sans dictionnaire.\cite{II-3-2-ref54}



\subsection{La lecture "étroite" (\textit{narrow reading})}

Pour pallier l'immensité du lexique grec, la stratégie de \textit{Narrow Reading} (lecture ciblée)
est particulièrement efficace. Elle consiste à lire exclusivement un auteur ou un genre spécifique
pendant une période donnée (par exemple, uniquement les Évangiles, ou uniquement Xénophon).

\begin{itemize}\item Cette approche réduit la charge lexicale globale, car chaque auteur a un idiolecte et un champ lexical privilégiés qui se répètent.\item La répétition naturelle du vocabulaire spécifique à l'auteur permet d'atteindre plus vite le seuil d'automaticité pour ce corpus précis.\cite{II-3-2-ref57}\end{itemize}

\subsection{L'apport des méthodes orales et communicatives (ci/tprs)}

L'automatisation lexicale (passage du déclaratif au procédural) est grandement facilitée par l'usage
oral de la langue, comme le pratiquent l'Institut Polis (Christophe Rico) ou les enseignants
utilisant le TPRS (\textit{Teaching Proficiency through Reading and Storytelling}).

\begin{itemize}\item \textbf{Living Sequential Expression (LSE) :} L'Institut Polis utilise cette technique inspirée de Gouin, où l'élève effectue et décrit des séries d'actions (ex: \og Je me lève, je marche vers la porte, j'ouvre la porte\fg{}). L'association du mouvement physique, de l'oralisation et de l'écoute crée une trace mnésique multimodale forte qui ancre le vocabulaire plus profondément que la simple lecture visuelle.\cite{II-3-2-ref58}\item \textbf{Pression Temporelle de l'Oral :} L'interaction orale impose une vitesse de traitement qui interdit le recours à la traduction consciente ou au dictionnaire mental. Elle force le cerveau à développer des accès directs (ventraux) au sens, accélérant ainsi l'automatisation nécessaire à la lecture fluide.\cite{II-3-2-ref40}\end{itemize}

\subsection{Enseignement explicite vs implicite du vocabulaire}

Bien que l'acquisition incidentelle (par la lecture) soit idéale pour la profondeur de la
connaissance, la recherche montre que pour atteindre les premiers paliers (les 2 000 mots les plus
fréquents), un enseignement explicite couplé à des systèmes de répétition espacée (SRS \- type Anki)
est un complément nécessaire pour accélérer la reconnaissance initiale. Cependant, cet apprentissage
explicite doit immédiatement être \og activé\fg{} par des lectures comprenant ces mots, sinon il
reste une connaissance inerte (déclarative) inutile pour la lecture fluide.\cite{II-3-2-ref38}

L'analyse de la littérature scientifique concernant le point II.3.\cite{II-3-2-ref2} de la thèse
confirme sans équivoque que le modèle traditionnel de lecture du grec ancien, basé sur le décryptage
grammatical et l'usage intensif du dictionnaire, est une impasse cognitive.

1. \textbf{Preuve Statistique :} Les recherches de Nation et Laufer établissent fermement qu'une
couverture lexicale de \textbf{98\%} est requise pour une lecture autonome et fluide. En dessous de
ce seuil, et particulièrement sous les 90-95\%, la compréhension se dégrade et la devinette
contextuelle devient impossible.

2. \textbf{Impasse Cognitive :} La consultation du dictionnaire durant la lecture impose une
\textbf{surcharge de la mémoire de travail} et un effet d'attention divisée qui brisent la
construction du sens (macro-structure). On ne peut pas \og lire\fg{} et \og chercher\fg{} en même
temps ; ce sont deux processus cognitifs concurrents.

3. \textbf{Nécessité de l'Automatisation :} La fluidité de lecture dépend du recyclage neuronal de
la zone \textbf{VWFA} et du passage des connaissances de la mémoire déclarative vers la mémoire
procédurale. Cela ne peut s'acquérir que par une exposition massive à un \textbf{input
compréhensible}.

4. \textbf{Réponse Pédagogique :} L'enseignement du grec doit opérer un changement de paradigme,
passant de l'analyse de textes trop difficiles (décodage) à la lecture massive de textes adaptés
(acquisition). L'utilisation de \textbf{lectures graduées}, de \textbf{textes enchâssés} (tiered
readings) et de méthodes orales actives n'est pas une \og dilution\fg{} de la rigueur classique,
mais la condition \textit{sine qua non} pour respecter l'architecture cognitive du cerveau apprenant
et permettre l'accès réel, et non simulé, aux textes de l'Antiquité.



\subsection{Tableaux récapitulatifs des données}



\subsection{Tableau 1 : relation entre couverture lexicale et expérience de lecture (hu \& nation, 2000\)}

\begin{table}[h!]\small\centering\begin{tabular}{| p{2cm} | p{2.2cm} | p{2.2cm} | p{2.2cm} |}\hline\textbf{Couverture Lexicale} & \textbf{Densité de Mots Inconnus} & \textbf{Expérience de Lecture / Compréhension} & \textbf{Faisabilité en Grec Ancien (Est.)} \\ \hline\textbf{100\%} & 0 & \textbf{Lecture Fluide}. Compréhension totale et automatique. & Expert lisant un texte simple. \\ \hline\textbf{98\%} & 1 mot sur 50 & \textbf{Lecture Autonome (Plaisir)}. Inférence contextuelle efficace. Dictionnaire inutile. & Niveau visé pour la fluidité. \\ \hline\textbf{95\%} & 1 mot sur 20 & \textbf{Lecture Instructionnelle}. Compréhension possible mais effort soutenu. Inférence partielle. & Seuil minimal académique (Laufer). \\ \hline\textbf{90\%} & 1 mot sur 10 & \textbf{Seuil de Frustration}. Compréhension globale difficile. Fil narratif souvent perdu. & Étudiant moyen face à un texte adapté. \\ \hline\textbf{\< 80\%} & \> 1 mot sur 5 & \textbf{Décodage Pur (Non-Lecture)}. Surcharge cognitive. Inférence impossible. Dépendance totale au dictionnaire. & Situation typique de l'étudiant face à un texte authentique. \\ \hline\end{tabular}\end{table}

\subsection{Tableau 2 : comparaison des processus cognitifs}

\begin{table}[h!]\small\centering\begin{tabular}{| p{2.3cm} | p{3.5cm} | p{3.5cm} |}\hline\textbf{Dimension} & \textbf{Lecture Fluide (Automatisée)} & \textbf{Lecture par Décodage (Grammaire-Traduction \+ Dictionnaire)} \\ \hline\textbf{Accès Lexical} & Voie Ventrale (Directe, \<200ms). Reconnaissance globale (VWFA). & Voie Dorsale (Indirecte, lente). Analyse graphème-phonème \+ Lemmatisation consciente. \\ \hline\textbf{Mémoire Sollicitée} & Mémoire Procédurale (Implicite). & Mémoire Déclarative (Explicite). \\ \hline\textbf{Charge Cognitive} & Faible (Ressources allouées au sens/macro-structure). & Extrême (Ressources allouées à la forme/micro-structure). \\ \hline\textbf{Rôle du Contexte} & Facilite la prédiction et l'inférence. & Inutilisable (trop de \og trous\fg{} dans le contexte). \\ \hline\textbf{Résultat} & Compréhension et Rétention du message. & Traduction (glose) souvent sans compréhension profonde. \\ \hline\end{tabular}\end{table}